\paragraph{新增面向学习的诊断。}
这些探测将奖励设计与可观测的选择和交互行为联系起来。原始候选池的公开奖励选择得到 0.20 的平均隐藏质量，高于均匀随机选择的期望值 0.075，增益集中在 A。五条新增前沿模型对照揭示了值得扩展的实验维度：GPT 在 B-L2-s0 的仅执行反馈条件成功，结构化/标量条件未完成；C-L2-s0 的固定与主动视角均成功，调用数分别为 10 和 26。这些探索性选例为更大规模的配对实验提供方向；一条未完成的 Claude C 续跑予以排除。

小模型新增八组等预算实验，覆盖 C、D、两阶段 A 和 E，比较独立采样与公开反馈修正，每阶段三个候选。格式适配器只提供可执行格式，不提供隐藏答案。30 个候选中 26 个进入原生评分，使 benchmark 能将接口可执行性与几何规划能力分开诊断。严格完成为 0/8，29 个候选的可接受奖励为零；这明确了分阶段课程和更密集进展反馈的实验切入点，同时保持最终验收标准一致。

另两组 D 单例直接反馈实际形状奖励，各复用一个起点并增加两次生成，分别采用配对实验的首个方案和模型自己的首个公开可执行方案。两组终点奖励与起点相同，尚未达到严格完成。另一个安全 D 候选获得 0.4 部分奖励，转运仍未完成，展示了可接受进展与端到端成功的区分。这些探测共同提供了可审计的课程学习与 RL 研究起点；扩大训练并开展留出集评估，将进一步检验其对策略提升的作用。
